\documentclass[11pt]{article}
\usepackage[margin=1in]{geometry}
\usepackage{amsmath,amssymb}
\usepackage{booktabs}
\usepackage{graphicx}
\usepackage{hyperref}
\usepackage{natbib}
\usepackage{xcolor}

\title{Latent Preference Inference via Conversational Active Learning:\\
  Cross-Domain Generalization Results}
\author{Patrick Cooper}
\date{\today}

\begin{document}
\maketitle

% ---------------------------------------------------------------------------
\section{Introduction}

This report presents the generalization study for the latent preference
inference system.  The system infers latent user preference parameters
(discount factor $\gamma$, risk aversion $\alpha$, loss aversion $\lambda$)
from structured binary-choice interactions, then transfers the inferred
parameters to produce personalized recommendations in new decision domains.

Three environments are evaluated:
\begin{enumerate}
  \item \textbf{Resource Strategy Game} -- a sequential allocation game with
    regime switching (4-channel variant A; 6-channel variant B).
  \item \textbf{Stock Backtesting} -- a multi-asset portfolio with five
    historical-calibrated asset classes.
  \item \textbf{Supply Chain Procurement} -- budget allocation across five
    suppliers with different cost, reliability, and disruption profiles.
\end{enumerate}

The study tests four hypotheses (README Section 8.3):
\begin{description}
  \item[H1] Within-domain personalization significantly outperforms generic
    (non-personalized) recommendations.
  \item[H2] Cross-domain transfer outperforms generic but underperforms
    within-domain personalization.
  \item[H3] The magnitude of cross-domain transfer correlates with structural
    similarity between domains; $\gamma$ (time preference) transfers better
    than $\alpha$ or $\lambda$ (domain-specific risk parameters).
  \item[H4] Active learning elicitation converges faster than random querying
    across all domains.
\end{description}

% ---------------------------------------------------------------------------
\section{Methods}

\subsection{Environments}

All environments extend a shared \texttt{BaseEnvironment} interface and expose
the same observation layout (wealth vector, market-state matrix, round index),
action space (simplex over $K$ channels), and preference-aware API
(\texttt{quality\_score}, \texttt{check\_quality\_floor},
\texttt{get\_optimal\_action}).

Regime switching uses a $3 \times 3$ Markov transition matrix shared across
all domains (bull/bear/neutral for the game and stock; normal/disrupted/crisis
for supply chain).

\subsection{Elicitation}

Preference elicitation uses a particle-filter posterior with EIG-maximizing
binary-choice queries.  Each domain has its own
\texttt{ScenarioLibraryBase} subclass generating $\gamma$, $\alpha$, and
$\lambda$-discriminating scenarios.  Gamma scenarios use multi-period
compounding horizons so that the discount factor is identifiable.

\subsection{Statistical Tests}

All pairwise comparisons use the Wilcoxon signed-rank test (one-sided where
directional hypotheses are stated).  Effect sizes are reported as rank-biserial
$r$ with Cohen's $d$ as a secondary measure.  Multiple comparisons across the
full family of tests are corrected with the Holm--Bonferroni step-down
procedure.

Theta stability is measured via ICC(2,1) (two-way random, single measures)
computed from an ANOVA decomposition.

% ---------------------------------------------------------------------------
\section{Experimental Setup}

\subsection{Domain Pairs}

Four directed domain pairs are evaluated:
\begin{itemize}
  \item Game variant A $\to$ Game variant B (structural similarity $\approx 0.95$)
  \item Game variant A $\to$ Stock ($\approx 0.55$)
  \item Game variant A $\to$ Supply Chain ($\approx 0.40$)
  \item Stock $\to$ Supply Chain ($\approx 0.50$)
\end{itemize}

\subsection{Parameters}

% Fill in after large-scale run
\begin{table}[h]
\centering
\begin{tabular}{ll}
\toprule
Parameter & Value \\
\midrule
Synthetic users per condition & \textit{(n\_users)} \\
Elicitation rounds & \textit{(max\_rounds)} \\
Particles & \textit{(n\_particles)} \\
EIG samples & \textit{(n\_eig\_samples)} \\
Stability sessions & 2 \\
\bottomrule
\end{tabular}
\caption{Experimental parameters.}
\label{tab:params}
\end{table}

% ---------------------------------------------------------------------------
\section{Results}

\subsection{H1: Within-Domain Personalization}

% Placeholder for H1 bar chart
% \begin{figure}[h]
%   \centering
%   \includegraphics[width=0.8\textwidth]{figures/hypothesis_panel.png}
%   \caption{H1--H4 hypothesis test results.}
%   \label{fig:hypothesis_panel}
% \end{figure}

\textit{Results to be populated from \texttt{outputs/generalization/results\_bundle.json}.}

\subsection{H2: Cross-Domain Transfer}

\textit{Cross-domain alignment scores vs.\ generic and within-domain baselines,
per domain pair.}

\subsection{H3: Parameter-Level Transfer}

% Placeholder for heatmap
% \begin{figure}[h]
%   \centering
%   \includegraphics[width=0.7\textwidth]{figures/parameter_heatmap.png}
%   \caption{Per-parameter cross-domain MAE.}
%   \label{fig:param_heatmap}
% \end{figure}

\textit{Per-parameter MAE decomposition showing whether $\gamma$ transfers
better than $\alpha$ or $\lambda$.}

\subsection{H4: Active vs.\ Random Elicitation}

\textit{Error reduction percentages and paired test results across all three
domains.}

\subsection{Theta Stability}

\textit{ICC(2,1) and Pearson correlations for test-retest reliability, per
parameter, per domain.}

% ---------------------------------------------------------------------------
\section{Discussion}

\subsection{Limitations}

\begin{itemize}
  \item The number of directed domain pairs (four) limits the power of the
    H3 similarity--transfer correlation; bootstrap confidence intervals
    supplement the Spearman test.
  \item All users are synthetic (softmax-rational).  Transfer to real human
    decision noise remains untested.
  \item Supply chain and stock environments use simplified regime models;
    real-world fidelity is limited.
\end{itemize}

\subsection{Structural Similarity Interpretation}

The structural similarity scores ($0.40$--$0.95$) are assigned based on
shared architecture (same $\gamma$/$\alpha$/$\lambda$ parameterization,
regime switching) and domain semantics.  Future work could learn similarity
from data.

% ---------------------------------------------------------------------------
\appendix
\section{Full Metric Tables}

\textit{Exported from \texttt{outputs/generalization/results\_bundle.json}.}

% Placeholder tables
% \input{tables/h1_results.tex}
% \input{tables/h2_results.tex}
% \input{tables/transfer_matrix.tex}
% \input{tables/stability.tex}

\end{document}
